\section{Linear and unitary components}\label{sec:linear}

This section applies the frame construction to linear and unitary groups.
At two we construct a coefficient to which Lemma~\ref{lem:propagation}
applies.  At odd primes we need the stronger conclusion that the coefficient
lies in the largest possible degree.  Projective commutation is the main
issue in both arguments: commuting projective transformations need not have
commuting matrix representatives.
The dimension-four cases at two are treated through their orthogonal
realizations in Propositions~\ref{prop:orthogonal} and~\ref{prop:psl45}.

\subsection{Weights, signs, and projective commutators}

We first record two elementary observations about weights.  If
$\Lambda$ is a finite subset of an affine space over $\F_p$, and translation
by a nonzero vector $a$ preserves $\Lambda$, then
\begin{equation}\label{eq:linear-translation}
 \dim\operatorname{aff}(\Lambda)\le \frac{|\Lambda|}{p}.
\end{equation}
Indeed, the translation has orbits of length $p$.  If there are $t$ orbits,
choose representatives $\lambda_1,\ldots,\lambda_t$; the affine span is
generated by $a$ and the $t-1$ differences $\lambda_i-\lambda_1$.

For the second observation, let an elementary abelian $p$-group $D$ act
faithfully projectively on a vector space in characteristic different from
$p$.  Suppose that its matrix representatives commute.  Over an algebraic
closure, rescale representatives of a basis of $D$ to have order $p$.
They give a linear representation of $D$ with a set of characters
$\Lambda\subseteq D^*$.  The differences of these characters span $D^*$:
an element annihilated by every difference acts by a scalar and is therefore
trivial in $D$.  Thus
\begin{equation}\label{eq:linear-affine-rank}
 \dim\operatorname{aff}(\Lambda)=\rk D.
\end{equation}
A transformation centralizing the projective action, and acting trivially
on the $p$th roots of unity, permutes these characters by a common
translation.  To see this, conjugating each chosen representative changes
it by a scalar; the scalars, which are $p$th roots of unity after
normalization, form a character of $D$.

Here is the complementary estimate when the representatives do not commute.
It will also be useful in the exceptional case of Section~\ref{sec:exceptional}.

\begin{lemma}\label{lem:projective-rank}
Let $D$ be an elementary abelian $p$-subgroup of $\PGL(V)$, where
$\dim V=n$ and the characteristic of the field is different from $p$.
The scalar commutators of representatives define an alternating bilinear
pairing on $D$ with values in the $p$th roots of unity.  If this pairing
has rank $2a$, then
\[
 p^a\mid n,
 \qquad
 \rk D\le 2a+\frac{n}{p^a}-1.
\]
In particular, for odd $p$, $n\ge5$, and $a>0$,
\begin{equation}\label{eq:odd-pairing-strict}
 \rk D<n-2.
\end{equation}
\end{lemma}

\begin{proof}
We may extend the field to its algebraic closure.  If $X,Y$ represent
elements of $D$, then $[X,Y]$ is scalar, and $X^p$ is scalar.  Consequently
$[X,Y]^p=[X^p,Y]=1$.  The commutator identities for central commutators
give the asserted alternating bilinear pairing.

Let $R$ be its radical.  Choose representatives of generators of $R$ and
rescale them to order $p$.  They commute with all the representatives in
question.  Let $t$ be the number of their distinct simultaneous weights.
Their weight ratios separate $R$ projectively, so $\rk R\le t-1$.
Choose a symplectic basis $x_1,y_1,\ldots,x_a,y_a$ for a complement of $R$
in $D$, with representatives normalized to order $p$.  On each radical
weight space, the commuting $x_i$ have simultaneous eigenspaces.
The $y_i$ translate their $a$ eigenvalue coordinates independently through
all $p$ possibilities.  They therefore permute these eigenspaces freely
in orbits of size $p^a$, with equal dimensions within an orbit.
Every radical weight space has positive dimension divisible by $p^a$.
It follows that $p^a\mid n$, $t\le n/p^a$, and
$\rk D=2a+\rk R\le 2a+n/p^a-1$.

For the final assertion put $u=p^a$.  Except when $(p,a)=(3,1)$,
we have $u>2a+2$: for $a=1$ this follows from $p\ge5$, and for $a\ge2$
it follows from $3^a>2a+2$.  Since $n\ge u$,
\[
 n-2-\left(2a+\frac n u-1\right)
 =n\left(1-\frac1u\right)-2a-1
 \ge u-2a-2>0.
\]
In the remaining case $3\mid n$ and $n\ge5$, so $n\ge6$, and the
same difference is $2n/3-3\ge1$.
\end{proof}

Let $L=\PSL_n(q)$ or $\PSU_n(q)$, and put $M=\PGL_n(q)$ or
$\PGU_n(q)$, respectively.  In the unitary notation, $\GU_n(q)$ denotes
the group of isometries of a nondegenerate hermitian form over
$\F_{q^2}$; its projective image also contains every projective
similitude, since the norm map onto $\F_q^\times$ is surjective.
In the linear case set $C=\F_q^\times$, and in the unitary case set
$C=\mu_{q+1}\subseteq\F_{q^2}^\times$.  Determinant induces
\begin{equation}\label{eq:linear-determinant}
 M/L\cong C/C^n.
\end{equation}
In particular this quotient is cyclic.

A line frame $\mathcal D$ is an unordered direct-sum decomposition into
$n$ lines; in the unitary case the lines are required to be orthogonal
and nondegenerate.  For a prime $p\mid |C|$, define $S(\mathcal D)$ to
be the projective images of the diagonal operators with entries in
$\mu_p$.  Thus $S(\mathcal D)\cong\F_p^n/\langle(1,\ldots,1)\rangle$
has rank $n-1$.  For any fixed $L\le H\le\Aut(L)$,
\begin{equation}\label{eq:linear-hzero}
 h_0:=\rk(S(\mathcal D)\cap H)\in\{n-2,n-1\}
\end{equation}
is independent of $\mathcal D$.  Indeed all these frames are conjugate
under $M$, and their images in the abelian group $M/L$ are consequently
the same subgroup, of order at most $p$.

We use the standard automorphism description for these groups in
dimension $n\ge5$; see \citet[Section~2 and statements~3.2--3.6]{Steinberg1960}
and \citet{GLS1998}.  If $q=r^f$, the quotient
$\Aut(\PSL_n(q))/\PGL_n(q)$ is the product of the field group of order
$f$ and the diagram group of order two.  The diagram operation is
$\gamma:g\mapsto(g^{-1})^{\mathsf t}$.  The quotient
$\Aut(\PSU_n(q))/\PGU_n(q)$ is the cyclic field group of order $2f$
on $\F_{q^2}$.  Semilinear isometries represent the latter automorphisms.
Thus an elementary subgroup has at most two external directions at two
in the linear case, at most one in the unitary case, and at most one
at odd primes in either case.

A \emph{correlation} here is a linear or semilinear identification
with the dual space, acting on the group through inverse transpose.
A dual frame consists of the lines of linear functionals vanishing
on all but one summand of the original frame.
A frame is matched with its dual frame when each of its lines is
matched with the corresponding dual line under this identification.

\begin{lemma}\label{lem:linear-sign-centralizer}
Suppose $q$ is odd and $n\ge5$.  Set
$S=S(\mathcal D)$ and $S_L=S\cap L$ for a line frame as above, using
$p=2$.  Every automorphism centralizing $S_L$ preserves each line of
$\mathcal D$, or matches it with the corresponding line of the dual
frame when it is a correlation.  It centralizes $S$.  Every involution
in $M$ centralizing $S_L$ belongs to $S$.
\end{lemma}

\begin{proof}
The preimage of $S_L$ in $\F_2^n$ is either all of $\F_2^n$ or
the kernel of the total-parity functional.  This follows directly from
the determinant formula; when $n$ is odd the second possibility cannot
occur, since the simultaneous sign belongs to that preimage.
For $n\ge5$, the coordinate characters on this preimage are distinct:
a relation equating two coordinates is neither the zero relation nor
the total-parity relation.  Their differences factor through $S_L$ and
span its dual.  After choosing a normalization, the corresponding $n$
weights therefore have affine rank $\rk S_L\ge n-2$.

Field operations and inverse-transpose leave binary sign values unchanged.
Projective centralization can consequently permute the weights only by a
common translation.  A nonzero translation would give affine rank at most
$n/2$ by \eqref{eq:linear-translation}, contrary to $n-2>n/2$.
Every weight line is therefore preserved individually, with the indicated
dual interpretation for a correlation.  Each such operation centralizes
all binary diagonal signs.  Finally a line-preserving projective
involution in $M$ has diagonal entries whose ratios square to one;
after removing a common scalar these entries are all $1$ or $-1$.
It belongs to $S$.
\end{proof}

\subsection{Compatible frames for involutory decorations}
\label{subsec:linear-compatible}

We next identify the frame spaces to which Theorem~\ref{thm:frame-bound}
will be applied.  A decoration is an involution outside $M$.  A frame is
compatible with it if every line is stable, with matching to the dual
frame in the correlation case.  For several decorations we require them
to commute and require compatibility with each of them.  In what follows
we consider only tuples having at least one compatible frame.

We use the following fixed-vector form of Galois descent. Let $K/k$
be a cyclic finite-field extension of degree $d$, with generator
$\sigma$, and let $B$ be an additive map on a finite-dimensional
$K$-space $V$ satisfying $B(av)=\sigma(a)B(v)$ and $B^d=\id$.
The powers of $B$ define a semilinear Galois action, and
\citet[Proposition~5.5]{MilneDescent2024} gives
\[
 K\otimes_k V^B\longrightarrow V,\qquad a\otimes v\longmapsto av,
\]
as an isomorphism. In particular $\dim_k V^B=\dim_K V$.
Every $B$-stable $K$-subspace descends by the same assertion applied
to its restriction; see also \citet[Proposition~3.5]{MilneDescent2024}.
Thus a stable line has a fixed nonzero generator.

After a form $f$ has been made $\sigma$-equivariant, meaning
$f(Bv,Bw)=\sigma(f(v,w))$, its Gram matrix in a fixed-vector basis
has entries in $k$. Its restriction consequently
extends back to the original form and preserves nondegeneracy.
For a hermitian form with field involution $\tau$ commuting with
$\sigma$, the restricted form has involution $\tau|_k$; it is
symmetric when that restriction is the identity. All field
involutions used here commute. The normalizations giving $B^d=\id$
and equivariance of the form are checked separately below; scalar
Hilbert~90, as in \citet[Corollary~5.25]{MilneFields2022}, handles
only the scalar part of those checks.

In the linear case, an external field involution requires $q=s^2$;
write $\sigma:x\mapsto x^s$.  On a compatible frame a representative
is $T=\diag(a_1,\ldots,a_n)\sigma$, and $T^2=cI$ implies
$a_i a_i^\sigma=c$ for every $i$.  Multiplying $T$ by $a_1^{-1}$ makes
$T^2=I$.  Semilinear descent gives an $n$-dimensional fixed space
$V_0$ over $\F_s$, with $V=V_0\otimes_{\F_s}\F_q$.
A stable line has a fixed generator by the one-dimensional norm-one
equation, so compatible frames correspond exactly to line frames of
$V_0$.

A correlation without field action is represented by a nonsingular
bilinear form.  The involutory condition says that its transpose is a
scalar multiple of itself.  In coordinates from a compatible frame its
matrix is diagonal and nonsingular; the scalar is therefore one.
The form is symmetric, and the compatible frames are its orthogonal
nondegenerate line frames.  For a diagram-field involution, the same
argument uses a $\sigma$-sesquilinear form.  Its conjugate transpose is
a scalar multiple of itself; scaling the form, using the norm-one scalar
equation, makes it hermitian.  Its compatible frames are hermitian frames
with parameter $s$.  Alternating forms have no compatible nondegenerate
line frame and do not occur in either assertion.

If there are two independent external involutions, choose their external
basis to be field and diagram-field.  Descend the first one to $\sigma$.
Choose generators for the lines of one compatible frame in the descended
space.  In these coordinates write the other involution as
$\operatorname{Int}(D)\gamma\sigma$, where $D$ is nonsingular diagonal
and $\operatorname{Int}(D)(g)=DgD^{-1}$.  Projective
commutation with $\sigma$ gives $D^\sigma=cD$.  Dividing by any diagonal
entry makes every entry of $D$ lie in $\F_s$.  The corresponding
correlation form on $V_0$ has Gram matrix $D^{-1}$.  Indeed, for
$g\in\GL(V_0)$ and $\lambda\in\F_s^\times$,
\[
 gDg^{\mathsf T}=\lambda D
 \quad\Longleftrightarrow\quad
 g^{\mathsf T}D^{-1}g=\lambda D^{-1}.
\]
This form is symmetric and nondegenerate over $\F_s$.
A frame is compatible with both
involutions precisely when it is the scalar extension of an orthogonal
frame for this symmetric form: the first condition descends its lines,
and the second is orthogonality for that form.

In the unitary case the unique external involution induces
$\tau:x\mapsto x^q$ on $\F_{q^2}$.  Take an orthonormal basis along a
compatible frame.  A unitary semilinear representative has the form
\[
 T=\diag(t_1,\ldots,t_n)\tau,
 \qquad t_i^{q+1}=1.
\]
It follows that $T^2=I$.  This conclusion uses compatibility with a
\emph{nondegenerate orthogonal} frame.  An arbitrary projective unitary
semilinear involution can instead have square $-I$, which cannot be
removed by multiplying by a norm-one scalar.  No such involution is
included here.  The fixed space of $T$ has dimension $n$ over $\F_q$.
The original hermitian form restricts to a nondegenerate symmetric form
on this space: its values on fixed vectors are fixed by $\tau$, and
hermitian symmetry then becomes ordinary symmetry.  Compatible unitary
frames correspond exactly to orthogonal nondegenerate line frames in
this fixed space.

These descriptions also identify all coarsenings of frames.  A sum of
lines descends to the corresponding sum of fixed lines, and the subgroup
requiring equal signs on that sum is unchanged by the description.
Adjoining the fixed decorations preserves these identifications.  Thus
the frame relations in Theorem~\ref{thm:frame-bound} are actual subgroup
face relations for the decorated groups.

\subsection{The construction at two}

\begin{proposition}\label{prop:linear-two}
Let $G$ have least order among finite groups with $O_2(G)=1$ and
$\widetilde H_*(\mathcal A_2(G);\Q)=0$.
Then $G$ has no simple component isomorphic to $\PSL_n(q)$ or
$\PSU_n(q)$ with $n\ge5$ and $q$ of characteristic at least five.
\end{proposition}

\begin{proof}
Suppose $L$ is such a component.  Consider all elementary abelian
$2$-subgroups $A\le N_G(L)$ for which $LA$ acts faithfully on $L$,
$O_2(C_G(LA))=1$, and $A\cap L$ contains $S(\mathcal D)\cap L$
for some line frame $\mathcal D$.  This family is nonempty: take
$A=S(\mathcal D)\cap L$ and use Lemma~\ref{lem:component-core}.
It is closed under elementary overgroups that are faithful
configurations, since each such overgroup still contains the same
subgroup $S(\mathcal D)\cap L$.
Choose a member of maximum rank, put $H=LA$, and identify $H$ with its
faithful image in $\Aut(L)$.

For its frame put $S=S(\mathcal D)$.  Lemma~\ref{lem:linear-sign-centralizer}
shows that $A$ centralizes $S$ and $A\cap M\le S$.  Adjoining $S\cap H$
would give another member of the same family with the same generated
group $H$, so maximality implies
\[
 A\cap M=S\cap H.
\]
Let $J_0$ be the image of $A$ in $\Aut(L)/M$ and write $d=\rk J_0$.
Choose a fixed basis $j_1,\ldots,j_d$ of $J_0$, and choose its lifts
$b_1,\ldots,b_d$ in a complement $T$ to $A\cap M$ in $A$.
For $d=2$ use the field and diagram-field basis above.  These lifts are
commuting involutions compatible with $\mathcal D$.  Filling out the
signs gives
\begin{equation}\label{eq:linear-E-intersection}
 E=S\times T,
 \qquad \rk E=n-1+d,
 \qquad E\cap H=A.
\end{equation}

For the construction, vary the frame and the ordered decorations subject
to the following requirements: $b_i\in H$ maps to $j_i$, the $b_i$ are
commuting involutions, and the frame is compatible with all of them.
The original data show that this set is nonempty.  All the resulting
groups $E$ have rank $n-1+d$.  They also have the same intersection index
in $H$, and every intersection generates $H$ over $L$.
Here is the verification of the latter point.  The image
$B$ of $S(\mathcal D)$ in $M/L$ is independent of the frame.
For the original $A$, its subgroup $A\cap M$ maps onto
$(H\cap M)/L$, since $A$ maps onto $H/L$.  Therefore $B$ contains
$(H\cap M)/L$.  For every new frame, $S(\mathcal D)\cap H$ maps onto
that same subgroup.  Together with the $b_i$ it consequently generates
$H/L$.  If $h_0$ is as in \eqref{eq:linear-hzero}, then for every datum
\begin{equation}\label{eq:linear-constant-index}
\begin{gathered}
 A'=E\cap H=(S(\mathcal D)\cap H)\times T,
 \qquad LA'=H,\\
 \rk A'=h_0+d,\qquad |E:A'|=2^c,\qquad
 c=n-1-h_0\in\{0,1\}.
\end{gathered}
\end{equation}

For each fixed decoration tuple, the preceding descent gives all the
frames of a linear, symmetric, or hermitian space of dimension $n$.
Its field parameter is $q$ or a subfield parameter $s$, and in every
case is at least five.  In the symmetric cases all line norm types are
allowed; the estimate is uniform over the isometry types that arise.
By Theorem~\ref{thm:frame-bound}, its space of
frame assignments has dimension at least the fraction
\[
 F=\left(\frac{17}{25}\right)^n
\]
of the number of compatible frames times $(n-1)!$.
We now count completions after one decoration has been omitted.

Fix a frame and one existing completion.  Multiplying the omitted
decoration by a diagonal element $X$ preserves compatibility.  The
following cyclic choices for each entry of $X$ also preserve its order
and its commutation with all retained decorations:
\[
\begin{array}{c|c|c}
\text{omitted direction}&\text{condition on an entry }x&k\\
\hline
\text{linear diagram}&x\in\F_q^\times&q-1\\
\text{linear field},\ q=s^2&x^{s+1}=1&s+1\\
\text{linear diagram-field},\ q=s^2&x\in\F_s^\times&s-1\\
\text{unitary external involution}&x^{q+1}=1&q+1
\end{array}
\]
In the two-direction case the last two linear prescriptions apply.
Indeed the field direction acts on diagonal entries by $x\mapsto x^s$,
and the diagram-field direction by $x\mapsto x^{-s}$.  Thus a norm-one
modification of the field lift is fixed by the retained diagram-field
lift, whereas a fixed-field modification of the diagram-field lift is
fixed by the retained field lift.  The same formulas show that the
square of the modified lift remains projectively trivial.  They apply
even when the original matrix representatives commute only projectively.

There are $k^n$ diagonal choices and $k^{n-1}$ projective choices.
We keep those projective modifications that lie in $L$; this ensures
that the new decoration remains in $H$ with the required external image.
The determinant homomorphism from the modification group to the cyclic
group $M/L$ has image of exponent dividing $k$, hence order at most $k$.
There are therefore at least $k^{n-2}$ permitted modifications.
Distinct projective modifications give distinct decorations.  Since
$k\ge4$, every partial datum has at least $4^{n-2}$ completions.
The inequality required in Lemma~\ref{lem:decorations} follows from
\begin{equation}\label{eq:linear-completion-inequality}
 \left(\frac{17}{25}\right)^n>\frac{2}{4^{n-2}}
 \qquad(n\ge5).
\end{equation}
At $n=5$ the left side exceeds $(2/3)^5=32/243>1/32$, which is the
right side; increasing $n$ multiplies their ratio by $68/25>1$.
For $d=0$ there are no decoration equations, and for $d=1$ the right
side can of course be halved.

It remains to verify that the assignments counted here are detected by
actual full-flag coefficients, as required by Lemma~\ref{lem:decorations}.
The top group $E$ recovers its full signs as $E\cap M$.
The simultaneous eigenspaces of these signs recover the frame.
For fixed $E$ and frame, take split flags through the full decoration
span $T$ and then through the sign flag.  If a second decoration span
$T'$ contributes to such a split flag, the vertex $T$ must split with
respect to $S\oplus T'$.  Its projection onto $J_0$ is an isomorphism
and its rank is $d$.  Hence its sign part is zero and $T=T'$.
The fixed basis $j_1,\ldots,j_d$ then recovers the individual $b_i$.
The coefficients above this vertex are precisely the original sign
coefficients, up to the fixed shuffle signs.  Thus these flags detect
all the counted parameters.  When $d=0$, detection is the undecorated
assertion of Theorem~\ref{thm:frame-bound}.

Lemma~\ref{lem:decorations} now supplies a nonzero homogeneous cycle
on full flags with top rank $n-1+d$ in the ambient automorphism group.
By \eqref{eq:linear-constant-index}, Lemma~\ref{lem:restriction} applies
at any supported top group: $c\le1<n-1+d$.  It produces a cycle in
$\mathcal A_2(H)$ with a nonzero coefficient on a saturated flag ending
at $A'=E\cap H$, of rank $h_0+d$.

We finally apply Corollary~\ref{cor:maximal-family} at this actual
coefficient.  The subgroup
$A'$ contains $S(\mathcal D)\cap L$, has the original maximum rank,
and satisfies $LA'=H$; thus $O_2(C_G(LA'))=O_2(C_G(H))=1$.
It therefore belongs to the family that was maximized.  The corollary
uses the nonzero saturated coefficient just constructed to give
$\widetilde H_*(\mathcal A_2(G);\Q)\ne0$, the desired contradiction.
\end{proof}

\subsection{Odd-prime unitary extensions}

\begin{proposition}\label{prop:unitary-odd}
Let $p$ be an odd prime, let $q$ be odd with $p\mid q+1$, and let
$L=\PSU_n(q)$ with $n\ge5$.  If $L\le H\le\Aut(L)$ and $H/L$ is
elementary abelian of $p$-power order, then
\[
 \widetilde H_{m_p(H)-1}(\Ap(H);\Q)\ne0.
\]
Here $m_p(H)$ denotes the maximum rank of an elementary abelian
$p$-subgroup of $H$.
\end{proposition}

\begin{proof}
Put $K=\F_{q^2}$ and $M=\PGU_n(q)$, and use the $p$-sign group of
an orthogonal line frame.  Its intersection with $H$ has the constant
rank $h_0\in\{n-2,n-1\}$ of \eqref{eq:linear-hzero}.
We first prove that
\begin{equation}\label{eq:unitary-rank-dichotomy}
 m_p(H)\in\{h_0,h_0+1\},
\end{equation}
and that in the second case the larger rank is attained by
$(S(\mathcal D)\cap H)\times\langle b\rangle$, where $b$ is a
field decoration preserving each line of $\mathcal D$.

Let $B\le H$ be elementary with $\rk B>h_0$, and set $D=B\cap M$.
There is at most one external direction, so $\rk D\ge n-2$.
By Lemma~\ref{lem:projective-rank}, the scalar commutator pairing of
$D$ is zero.  Choose commuting unitary representatives $U_1,\ldots,U_r$
of a basis of $D$, where $r=\rk D$, and write $U_i^p=c_iI$.
Over an algebraic closure choose $t_i^p=c_i$ and put $X_i=t_i^{-1}U_i$.
Fix a primitive $p$th root $\zeta$.  The simultaneous weights of the
$X_i$ form a set $\Lambda\subseteq\F_p^r$ with
$\dim\operatorname{aff}(\Lambda)=r$ and $|\Lambda|\le n$; the actual
eigenvalue of $U_i$ at $\lambda$ is $t_i\zeta^{\lambda_i}$.

Unitarity identifies this representation with its conjugate dual, so
the multiset of joint eigenvalues is preserved by raising every entry
to the power $-q$.  Since $-q\equiv1\pmod p$, this operation acts on
$\Lambda$ by the common translation $\lambda\mapsto\lambda+\delta$,
where
\[
 \zeta^{\delta_i}=t_i^{-(q+1)}.
\]
The right side is a $p$th root of unity because $c_i^{q+1}=1$.
A nonzero translation would imply
\[
 \rk D\le |\Lambda|/p\le n/p<n-2,
\]
by \eqref{eq:linear-translation}, which is impossible.  The translation
is therefore zero.  Hence $t_i^{q+1}=1$ for every $i$.  In particular
$t_i\in\mu_{q+1}\subseteq K$, so the $X_i$ are unitary matrices over
$K$ with eigenvalues in $\mu_p$.  Their simultaneous eigenspaces are
defined over $K$.  Distinct weights are orthogonal: if their $i$th
eigenvalues differ, invariance of the hermitian form multiplies their
pairing by the nonidentity ratio of those eigenvalues.  Each eigenspace
is nondegenerate since their orthogonal direct sum is the whole space.
They therefore admit orthogonal nondegenerate line refinements over $K$.
On such a refinement the $X_i$ are diagonal $p$-signs, so their
projective images belong to $S(\mathcal D)$.
Thus $D\le S(\mathcal D)\cap H$ for some frame.

If $B=D$, this contradicts $\rk B>h_0$.  Otherwise choose
$b\in B\setminus M$.  Its field action has order $p$ and hence fixes
the cyclic group $\mu_p$, whose automorphism group has order $p-1$.
Projective commutation with $D$ once more permutes its weights by a
translation.  The same rank inequality excludes a nonzero translation,
so $b$ preserves every nondegenerate eigenspace just obtained.

Write $q=s^p$ and let $\sigma$ be the order-$p$ field action of $b$
on $K$; its fixed field is $K_0=\F_{s^2}$.  A unitary semilinear
representative $T$ of $b$ satisfies $T^p=cI$, where
$c\in\mu_{q+1}$ and $c^\sigma=c$, so $c\in\mu_{s+1}$.
The norm
\begin{equation}\label{eq:unitary-norm}
 N:\mu_{q+1}\longrightarrow\mu_{s+1}
\end{equation}
for $K/K_0$ is onto.  Indeed, writing
\[
 k=\frac{s^p+1}{s+1},\qquad
 \ell=\frac{s^p-1}{s-1},
\]
its exponent is $(s^{2p}-1)/(s^2-1)=k\ell$.
Since $p$ is odd, $\ell=1+s+\cdots+s^{p-1}\equiv1\pmod{s+1}$.
As $|\mu_{q+1}|=k(s+1)$, the image has order $s+1$ and the kernel
has order $k$.  This norm is unchanged if another generator of
$\Gal(K/K_0)$ is used.  Choose $a\in\mu_{q+1}$ with $N(a)=c^{-1}$.
Replacing $T$ by $aT$ gives $T^p=I$ while preserving its unitary
semilinear action.

Semilinear descent now gives an $n$-dimensional space $V_0$ over $K_0$.
The hermitian form descends: on fixed vectors its values lie in $K_0$,
and the involution $x\mapsto x^q$ restricts there to $x\mapsto x^s$.
Each of the $T$-stable nondegenerate eigenspaces descends as well.
Choose an orthogonal line frame within each descended eigenspace and
extend scalars back to $K$.  The resulting frame $\mathcal D$ is
compatible with $b$, contains $D$ in its signs, and has each line fixed
by $b$.  Since $\sigma$ fixes $\mu_p$, $b$ centralizes all of
$S(\mathcal D)$.  Hence
\[
 B\le (S(\mathcal D)\cap H)\times\langle b\rangle,
\]
an elementary group of rank $h_0+1$.  This proves
\eqref{eq:unitary-rank-dichotomy} with the asserted compatible realization.

We now construct a cycle in the degree dictated by this calculation.
Set $c=n-1-h_0\in\{0,1\}$.  If $m_p(H)=h_0$, use
Theorem~\ref{thm:frame-bound} on all unitary frames with parameter $q$.
Here $q\ge5$: the smaller odd possibility $q=3$ has no odd prime
dividing $q+1$.  It gives a nonzero cycle with top rank $n-1$ in
$M$.  Its intersection index in $H$ is the constant $p^c$, so
Lemma~\ref{lem:restriction} produces a nonzero cycle in degree
\begin{equation}\label{eq:unitary-undecorated-degree}
 (n-1)-1-c=h_0-1=m_p(H)-1.
\end{equation}

Suppose instead that $m_p(H)=h_0+1$.  Fix the nontrivial field image
of a compatible element $b$ supplied above.  Use all pairs consisting
of an order-$p$ lift of this fixed image in $H$ and a compatible
orthogonal line frame.  For each lift that occurs, the normalization and
descent just proved identify its compatible frames with all hermitian
frames over the parameter $s=q^{1/p}$.  Fermat's congruence gives
$s\equiv s^p=q\equiv-1\pmod p$, so $p\mid s+1$.
As $s$ is odd, it follows that $s\ge5$.  The frame fraction is again
at least $(17/25)^n$.

At a fixed frame vary a given lift by a diagonal matrix with all entries
in the kernel of \eqref{eq:unitary-norm}.  The kernel has order
$k=(q+1)/(s+1)$, and the norm equation ensures that the modified lift
still has projective order $p$.  There are $k^{n-1}$ choices modulo
common scalars.  As in the determinant calculation at two, restricting
the modification to $L$ loses a factor at most $k$, leaving at least
$k^{n-2}$ completions in $H$.  To bound $k$, write $p=2a+1$ and note
\[
 k=1+\sum_{j=1}^{a}(s^{2j}-s^{2j-1})
 \ge s^2-s+1\ge21.
\]
Consequently
\[
 \left(\frac{17}{25}\right)^n>\frac1{21^{n-2}}
 \qquad(n\ge5).
\]
At $n=5$ the left side exceeds $1/32>1/9261$, and increasing $n$
multiplies the ratio of the two sides by $357/25$.
The detection argument used above applies without change: the full
group $E=S(\mathcal D)\times\langle b\rangle$ recovers its linear
kernel and frame, and the split flags through the line $\langle b\rangle$
recover the lift of the fixed field generator.  All the hypotheses of
Lemma~\ref{lem:decorations} are therefore satisfied.

We obtain a nonzero cycle with top rank $n$.  For every supporting group
$E$, the decoration lies in $H$ and
\[
 E\cap H=(S(\mathcal D)\cap H)\times\langle b\rangle,
 \qquad |E:E\cap H|=p^c.
\]
Lemma~\ref{lem:restriction} gives a nonzero cycle of degree
\begin{equation}\label{eq:unitary-decorated-degree}
 n-1-c=h_0=m_p(H)-1.
\end{equation}
In both cases the resulting cycle has a nonzero coefficient and lies
in the largest possible chain degree of $\Ap(H)$.  There are no chains
one degree higher, so this cycle cannot be a boundary.  This proves the
proposition.
\end{proof}
